\subsection*{導入：専門職実務が重要な理由}

\addcontentsline{toc}{subsection}{導入：専門職実務が重要な理由}

ソフトウェア工学専門職実務は、ソフトウェア技術者が専門的、責任ある、倫理的な方法で仕事をするために必要な知識、技能、態度を扱う知識領域である。専門職実務は、成果物だけでなく、作業の過程にも基準を求める。つまり「動くものを作ったからよい」ではなく、「どのような根拠、手順、責任、説明可能性をもって作ったか」も問われる。

専門職としての実践には、一般に受け入れられた知識体系、倫理規程と行動規範、認定・資格証明・免許の仕組み、専門職団体、標準が関わる。ソフトウェア技術者は、ACM、IEEE、IFIP、IEEEコンピュータソサエティ、ISO/IEC、ISO/IEC/IEEEなどが示す標準や指針を参考にしながら、自分の仕事を説明できる形で進める必要がある。

第14章は三つの大領域で構成される。第一は専門性であり、倫理、標準、法的論点、契約、文書化、トレードオフ分析を扱う。第二は集団力学と心理学であり、チーム作業、個人の認知、複雑な問題、利害関係者、不確実性、多様性を扱う。第三はコミュニケーション技能であり、読む、書く、チーム内で伝える、発表する能力を扱う。
